\section{REFERENCIAL TEÓRICO}
\label{sec:referencial}

Esta seção articula os estudos educacionais sobre os fatores associados ao desempenho, com foco no IDEB e no Brasil, e os métodos computacionais de mineração de dados, em especial o agrupamento como estratégia de segmentação contextual.

\subsection{O IDEB e os fatores associados ao desempenho escolar no Brasil}

O IDEB combina a taxa de aprovação (fluxo) e a proficiência em Língua Portuguesa e Matemática aferida pelo Sistema de Avaliação da Educação Básica (Saeb) \parencite{fernandes2007}. \textcite{soares2013} examinaram os pressupostos estatísticos do índice e alertaram para o risco de interpretações que ignoram a correlação entre fluxo e proficiência.

O debate sobre o efeito-escola remonta ao relatório de \textcite{coleman1966}, segundo o qual o contexto familiar explica a maior parte da variância no desempenho, e os fatores escolares têm efeito residual, mas acionável por políticas públicas. No Brasil, a pesquisa em eficácia escolar, sistematizada por \textcite{brooke2008}, mostrou que, controlado o nível socioeconômico dos alunos, as escolas ainda diferem no quanto seus estudantes aprendem. Como qualidade e equidade variam com o nível socioeconômico das escolas, comparações entre contextos distintos tendem a confundir efeito-escola com efeito-contexto \parencite{albernaz2002}.

Entre os fatores escolares, \textcite{rivkin2005} demonstraram que a qualidade do professor é o fator intraescolar de maior impacto sobre o aprendizado, mas que a titulação explica pouco dessa variação. A Meta 15 do Plano Nacional de Educação estabelece que todos os professores da educação básica tenham formação específica de nível superior, em licenciatura na área em que atuam \parencite{brasil2014}, e \textcite{gatti2019} descrevem um cenário de expansão das licenciaturas, sobretudo a distância, com persistência de docentes em exercício sem a formação exigida. Por isso, interessa verificar se a titulação se associa ao desempenho em redes nas quais parte expressiva dos docentes ainda não tem formação superior.

Quanto à infraestrutura, \textcite{alves2013} mostraram que o IDEB das escolas varia com o nível socioeconômico dos alunos, a infraestrutura e a complexidade da escola, com associação mais forte nos anos iniciais. \textcite{vasconcelos2021} encontraram, para os municípios brasileiros, que a infraestrutura pesa mais sobre o IDEB do que o volume de investimento em educação. Em revisão sobre países em desenvolvimento, \textcite{glewwe2013} verificaram que, nos estudos de maior rigor, poucas características de escolas e professores afetam o aprendizado, e a escolaridade docente não mostrou efeito consistente.

\subsection{Mineração de dados e segmentação contextual}

A mineração de dados educacionais (\textit{Educational Data Mining}) consolidou-se como área na segunda metade da década de 2000, com paradigmas que vão da predição ao agrupamento, mas é dominada por técnicas supervisionadas para predizer aprovação, evasão e desempenho \parencite{romero2020}. No Brasil, os microdados do Instituto Nacional de Estudos e Pesquisas Educacionais Anísio Teixeira (INEP) impulsionaram estudos que aplicam essas técnicas ao Censo Escolar e ao IDEB \parencite{colpani2018}. Na literatura consultada, porém, predominam modelos globais, que tratam como comparáveis municípios de contextos socioeconômicos muito distintos, e são menos frequentes os estudos sobre o Norte e o Nordeste.

O \textit{K-Means} \parencite{macqueen1967} agrupa entidades em grupos internamente homogêneos e exige definir previamente o número de grupos~$k$, que pode ser escolhido pelo \textit{Silhouette Score} \parencite{rousseeuw1987}: para cada ponto, ele compara a coesão com o próprio grupo e a separação em relação ao grupo mais próximo, num escore entre $-1$ e 1. No contexto brasileiro, \textcite{alves2014} usaram o \textit{K-Means} para agrupar as escolas de educação básica em sete níveis segundo o índice socioeconômico, o que mostra o potencial da técnica para caracterizar contextos antes de comparar os grupos em outros indicadores.

Combinar essa segmentação com testes aplicados dentro de cada grupo reduz, sem eliminar, o risco de atribuir a fatores escolares diferenças determinadas pela renda. É essa a lógica em dois estágios adotada neste estudo.
